\subsection{代数扭转与几何扭转}\label{sec:17j}
假设有一个以圆盘为底空间、纤维为 $M$、恰有一个临界点的 Lefschetz 纤维化。若纤维维数 $2n>2$，这样的纤维化总存在从 $M$ 的二次复体积形式继承而来的相对二次复体积形式；若 $2n=2$，则当且仅当消失圈 $V$ 能赋予分次时才存在。考虑到这一点，\secref{17d}至\secref{17h}中的全部论证都可照常进行，所得符号也都与\stref{引理}{5.3}中预期的符号一致；这一验证直接但略显繁琐，这里略去。如\stref{例题}{17.14}所示，当 $2n>2$ 时，标准边界条件不能赋予相对分次，但这不成问题，因为在这些维数中，\stref{笔记}{17.8}的消失论证不依赖任何抵消。在剩下的情形中，\stref{例题}{17.14}适用，因此若要在 $\operatorname{char}(\K)\ne2$ 的情形下工作，就必须在消失圈 $V=L_0$ 上选择非平凡 Pin 结构。\secref{17h}中的主要结果，在补上正确的分次和符号后，如下所述。

\StatementNumber{16}\begin{theorem}\label{st:17.16}
设 $V=L_0$ 为 $M$ 中带标架的正合 Lagrangian 球面，并带有膜结构 $L_0^\#$；当 $n=1$ 且 $\operatorname{char}(\K)\ne2$ 时，该结构中的 Pin 结构应非平凡。设 $L_1,L_2$ 为另外两个正合 Lagrangian 子流形，固定一个同痕 $L_2\simeq\tau_V(L_1)$，并假设它们也带有与该同痕相容的膜结构 $L_k^\#$，其中为 $\tau_V$ 赋予自然分次。以 $Y_k$ 表示 $L_k^\#$ 给出的 $\F(M)$ 的对象，以 $\Y_k$ 表示其 Yoneda $\Ai$ 模。那么，通过计数纤维化 $E^c\to S^c$ 与 $E^{k,r}\to S^{k,r}$ 的截面，可得到
\[
c\in CF^0(L_1^\#,L_2^\#),\qquad
k:CF^*(L_0^\#,L_1^\#)\longrightarrow CF^*(L_0^\#,L_2^\#)[-1],
\]
使按\stref{引理}{5.3}定义的相应 $[t]:T_{Y_0}(\Y_1)\to\Y_2$ 在 $H^0(\modcat(\F(M)))$ 中为同构。
\end{theorem}

为得到更具体的陈述，现在像\eref{5.6}中那样，以扭曲复形代替 $\Ai$ 模。再假设 $L_2=\tau_V(L_1)$，并略去一些细节，\stref{定理}{17.16}就成为下列“$T=\tau$”陈述：
\StatementNumber{17}\begin{corollary}\label{st:17.17}
在导出范畴 $D\F(M)$ 中，有 $T_{L_0^\#}(L_1^\#)\cong\tau_{L_0}(L_1^\#)$。
\end{corollary}
\SourcePage{270}
$T_{L_0^\#}(L_1^\#)$ 属于正合三角\eref{5.4}，由正合三角的一般性质，可将其换成 $\tau_{L_0}(L_1^\#)$。注意，三角中的一个态射为\secref{5b}所定义的纯代数求值态射 $[f]$。由\stref{引理}{5.3}，第二个态射为 $[c]$，因此有：
\clearpage
\enlargethispage{\baselineskip}
\StatementNumber{18}\begin{corollary}\label{st:17.18}
在 $D\F(M)$ 中存在正合三角
\SourceEquationDiagram{diagram-17-23}{0.413307247}{17.23}
\end{corollary}

\StatementNumber{19}\begin{note}\label{st:17.19}
\eref{17.23}中的第三个映射是下列空间中的元素：
\[
\begin{aligned}
&\Hom_{D\F(M)}\bigl(\tau_{L_0}(L_1^\#),HF^*(L_0^\#,L_1^\#)\otimes L_0^\#[1]\bigr)\\
&\quad\cong\bigl(HF^*(L_0^\#,L_1^\#)\otimes HF^*(\tau_{L_0}(L_1^\#),L_0^\#)\bigr)^1\\
&\quad\cong\bigl(HF^*(L_0^\#,L_1^\#)\otimes HF^*(L_1^\#,L_0^\#)\bigr)^n.
\end{aligned}
\]
最后一个同构使用了\stref{例题}{17.14}中提到的等式 $\tau_{L_0}(L_0^\#)=L_0^\#[1-n]$。\znote{262}{两个指数是同一笔平移账：目标中的 $[1]$ 使整个 $\Hom$ 落在 $\bigl(HF^*\otimes HF^*\bigr)^1$；又由 $\tau_{L_0}(L_0^\#)=L_0^\#[1-n]$，把 $\tau_{L_0}$ 同时作用到两个变元上得 $HF^*(\tau_{L_0}(L_1^\#),L_0^\#)\cong HF^*(L_1^\#,L_0^\#)[n-1]$，次数整体升 $n-1$，指数便由 $1$ 变成 $n$。这类次数账在正合三角里最容易错：只要把 $\tau_{L_0}(L_1^\#)$ 换成 $L_1^\#$ 而忘记 $L_0^\#$ 同时平移，指数就差 $n-1$。[来源：Seidel, Fukaya categories and Picard–Lefschetz theory, §17j 与例 17.14。]}自然会预期，这就是由一对 Poincaré 对偶基给出的元素；等价地，可以将其视为\eref{5.20}中所用的那种对偶求值映射。原则上，将\stref{笔记}{5.5}与\stref{笔记}{17.13}结合起来，似乎就能证明情况确实如此。不过，我们尚未核查全部细节；尤其是符号问题颇为棘手，会以非平凡的方式影响结果，参见\stref{笔记}{3.12}。
\end{note}

最后，给出一个更有分量的应用。设 $V^\#=(V_1^\#,\ldots,V_m^\#)$ 为 $M$ 中一组带标架的正合 Lagrangian 球面，每个都带有膜结构；当 $n=1$ 且 $\operatorname{char}(\K)\ne2$ 时，要求所有 Pin 结构非平凡。设 $W^\#=\LG_kV^\#$ 为对它作一次初等 Hurwitz 移动\eref{16.7}得到的集合，并赋予自然诱导的膜结构；其中的分次使用 $\tau_{V_k}$ 的自然分次。设 $\F^\to(V^\#)$ 与 $\F^\to(W^\#)$ 为 $\F(M)$ 中相应的有向 $\Ai$ 子范畴。\stref{推论}{17.17}表明，$W^\#$ 的各对象在 $D^b(\F(M))$ 中，与对 $V^\#$ 施以变异\eref{5.24}所得集合的各对象同构。由\stref{引理}{5.21}，换成同构对象不会改变有向 $\Ai$ 范畴的拟同构类。因此，\stref{引理}{5.23}适用，并说明 $\Tw\F^\to(V^\#)$ 与 $\Tw\F^\to(W^\#)$ 拟等价。另一个更初等得多的事实是：平移一个或多个 $V_j$ 的分次，也不会改变 $\Tw\F^\to(V^\#)$。将结论归纳如下：
\SourcePage{271}
\StatementNumber{20}\begin{theorem}\label{st:17.20}
在拟等价意义下，$\Tw\F^\to(V^\#)$ 是 $V$ 的 Hurwitz 等价类的不变量。
\end{theorem}
将其应用于消失圈的特选基，立即得到下述结论：
\StatementNumber{21}\begin{corollary}\label{st:17.21}
设 $\pi:E\to S$ 为底空间 $S$ 是圆盘的 Lefschetz 纤维化，并带有相对二次复体积形式 $\eta^2_{E/S}$。选取基点 $*$，为 $M=E_*$ 赋予相应的形式 $\eta_M^2$。设 $\symbfit{\gamma}$ 为从 $*$ 出发的消失路径特选基，以 $V_{\symbfit{\gamma}}$ 表示 $M$ 中相应的消失圈集合。为每个消失圈赋予分次和 Pin 结构；若 $n=1$ 且 $\operatorname{char}(\K)\ne2$，则选择非平凡 Pin 结构。将所得集合记为 $V_{\symbfit{\gamma}}^\#$，并设 $\F^\to(V_{\symbfit{\gamma}}^\#)$ 为 $\F(M)$ 中相应的有向 $\Ai$ 子范畴。那么，在拟等价意义下，扭曲复形范畴 $\Tw\F^\to(V_{\symbfit{\gamma}}^\#)$ 与 $\symbfit{\gamma}$ 及分次的选择无关，因此是该 Lefschetz 纤维化的不变量。故导出范畴 $D\F^\to(V_{\symbfit{\gamma}}^\#)$ 也具有同样的性质，只须将拟等价换成三角范畴的等价。
\end{corollary}
\StatementNumber{22}\begin{note}\label{st:17.22}
\stref{推论}{17.17}最初由 Kontsevich 猜想提出。显然，人们会期待一个更强的版本成立，即 $\tau$ 诱导的自同构，应在 $H^0(\operatorname{fun}(\F(M),\F(M)))$ 中典范地同构于 $T$。这里不证明这一点；不过，预期它将是 Wehrheim–Woodward 正在进行的 Lagrangian 对应研究的一个推论。
\end{note}
